\documentclass[../../main.tex]{subfiles}

\graphicspath{{../../image/}{../../figure/}}

\begin{document}

\section{计算机视觉}\label{section:ch10-计算机视觉}

图像数据的自动分析和解释是计算机视觉（computer vision）领域的焦点，代表了机器学习的一个主要应用领域（Szeliski, 2022）。计算机视觉主要基于三维投影几何，人们曾手工构建特征并用作简单学习算法的输入（Hartley and Zisserman, 2004）。然而，计算机视觉是最早被深度学习革命改变的领域之一，这主要归功于 CNN 架构。虽然该架构最初是在图像分析场景下开发的，但它也被应用于序列数据分析等其他场景。最近，基于 Transformer（参见\cref{chapter:ch12-Transformer与注意力机制}）的替代架构在一些应用中已经与 CNN 架构形成竞争态势。

机器学习在计算机视觉领域的应用有很多，其中最常见的一些应用如下：
\begin{enumerate}
  \item 图像分类，例如将一幅皮肤病变的图像分类为良性或恶性（参见\cref{figure:ch01-皮肤病变}），这有时称为“图像识别”；
  \item 目标检测，检测图像中的目标并确定它们在图像中的位置，例如从自动驾驶汽车收集的相机数据中检测行人（参见\cref{figure:ch10-边界框}）；
  \item 图像分割，其中每个像素被单独分类，从而将图像划分为共享相同标签的区域。例如，自然场景可能分割成天空、草地、树木和建筑物（参见\cref{figure:ch10-语义分割}），而医学扫描图像可以分割成癌变组织和正常组织；
  \item 标题生成，即从图像中自动生成文本描述（参见\cref{chapter:ch12-Transformer与注意力机制}）；
  \item 图像合成，例如生成人脸图像（参见\cref{figure:ch01-生成人脸}），也可以根据描述所需图像内容的文本输入来合成图像；
  \item 图像修补，将图像的一个区域替换为与图像其余部分一致的合成像素，这可用于在图像编辑过程中删除不需要的对象（参见\cref{chapter:ch20-扩散模型}）；
  \item 风格迁移，即将一种风格的输入图像（例如一张照片）转换为不同风格的相应图像，例如一幅油画（参见\cref{figure:ch10-风格迁移}）；
  \item 图像超分，通过增加像素数量以及与合成相关的高频信息来提高图像的分辨率（参见\cref{chapter:ch20-扩散模型}）；
  \item 深度预测，即使用一个或多个视图来预测目标图像中每个像素处场景与相机的距离；
  \item 场景重建，即使用一个或多个场景的二维图像来重建三维表示。
\end{enumerate}

\subsection{图像数据}\label{subsection:ch10-图像数据}

图像由像素的矩形数组组成，其中每个像素要么具有灰度强度，要么具有更常见的红、绿、蓝通道的三元组强度值。这些强度值是非负数，同时也有最大值，对应于相机或用于捕捉图像的其他硬件设备的限制。大多数情况下，我们将强度视为连续变量，但在实践中它们是用有限精度表示的，例如用 8 比特表示范围 $0\sim 255$ 内的整数。一些图像，例如医学诊断中使用的磁共振成像（Magnetic Resonance Imaging, MRI）扫描图像，是包含体素的三维网格。类似地，视频包含一系列二维图像，因此也可以视为三维结构，其中连续的帧随着时间堆叠在一起。

\begin{definition}[图像数据]\label{definition:ch10-图像数据}
称由像素（pixel）构成的矩形数组为\textbf{图像}，其中每个像素携带一个灰度强度值，或携带红（R）、绿（G）、蓝（B）三个通道的强度三元组。若图像有 $J$ 行、$K$ 列、$C$ 个通道，则可用尺寸为 $J\times K\times C$ 的张量描述它，其中 $C=1$ 对应灰度图像、$C=3$ 对应彩色图像。强度取非负值且有上界，实践中以有限精度表示，例如以 8 比特整数表示 $0\sim 255$ 内的值。若图像元素排列成三维网格而不是二维网格，则称其元素为\textbf{体素}（voxel）；视频是沿时间轴堆叠的二维图像序列，因而也构成三维结构。
\end{definition}

% SUPPLEMENT: parent=definition:ch10-图像数据; kind=motivation; location=紧随其后的 remark（重构新增）
\begin{remark}（重构新增）
这个定义把“图像”从直观的图片转成可直接送入网络的对象：一个三阶张量。定义中的三件事各有用处。第一，通道维 $C$ 与空间维 $J\times K$ 分离，使\cref{definition:ch10-多维卷积}可以把“在空间上共享权重、在通道上求和”写成四阶滤波器张量。第二，强度有界（如 $0\sim 255$）意味着输入落在紧集中，于是梯度下降与梯度符号法（\cref{definition:ch10-对抗攻击}）中的像素更新必须再投影回该集合，否则会越界。第三，“视为连续变量还是有限精度”只是建模约定：把强度视为连续变量是为了可微，存储上的离散化不改变前向计算。注意体素网格与视频序列虽然都是三维结构，但第三维的语义分别是空间深度与时间，滤波器在这两维上是否共享权重必须另行规定。
\end{remark}

考虑在图像数据上应用神经网络实现上述应用所面临的挑战。图像通常具有高维特性，通过相机捕获的图像包含数千万像素。因此，将图像数据视为非结构化数据可能需要一个具有大量参数的模型，而这会导致模型难以训练。更重要的是，这样的方法没有考虑到图像数据的高度结构化性质，其中不同像素的相对位置起着至关重要的作用。我们之所以关注这一点，是因为如果我们将图像的像素随机排列，那么结果就不再像自然图像了。如果我们为每个像素独立地随机采样像素强度来生成合成图像，那么生成的合成图像看起来像自然图像（参见\cref{figure:ch06-自然与随机图像}）的可能性基本为零。局部相关性很重要：在自然图像中，与两个相距很远的像素相比，两个邻近像素具有相似颜色和强度的概率要高得多。这表明我们可以将图像的先验知识通过强大的归纳偏置编码到神经网络中，从而得到参数少得多、泛化精度高得多的模型。

\end{document}
